\documentclass[11pt]{article}
\usepackage{times}
\usepackage[margin=1in]{geometry}
\usepackage{amsmath}
\usepackage{amssymb}
\usepackage{amsthm}
\usepackage{graphicx}
\usepackage{fancyhdr}
\usepackage{parskip}
\usepackage{sourcecodepro}
\usepackage{listings}
\usepackage{color}
\usepackage{hyperref}

% Header setup
\pagestyle{fancy}
\fancyhf{}
\rhead{Name: \underline{\hspace{3cm}}}
\lhead{APMA 4302 - Methods}
\cfoot{\thepage}

\begin{document}
\lstset{basicstyle=\ttfamily}

\title{APMA 4302 Methods - Homework 4: Solution of coupled multiphysics PDEs with Firedrake and PETSc}
\author{Marc Spiegelman}
\date{\today}
\maketitle

Here we will solve the coupled nonlinear system of dimensionless equations governing infinite prandtl number thermal convection in a 2D box, of height $H$which can be written as:
\begin{align}
  \frac{\partial T}{\partial t} + \mathbf{v} \cdot \nabla T & = \frac{1}{Ra}\mathbf{\nabla}^2 T  \label{eq:temp}\\
  -\nabla^2\mathbf{v} + \mathbf{\nabla} p & = T\hat{\mathbf{j}} \label{eq:stokes}\\
  \mathbf{\nabla} \cdot \mathbf{v} & =  0 \label{eq:incomp}
\end{align}
which couples advection diffusion of temperature, $T$, (Eq. \eqref{eq:temp}) with incompressible Stoke's flow for the velocity, $\mathbf{v}$ and pressure, $p$ (Eqs.\eqref{eq:stokes} and \eqref{eq:incomp}). Here $\hat{\mathbf{j}}$ is the unit vector in the vertical direction (pointing upward). When non-dimensionalized, the equations are  governed by a single dimensionless parameter, the Raleigh Number
\begin{displaymath}
  Ra = \frac{\rho_0\alpha g \Delta T H^3}{\kappa \eta },
\end{displaymath}
which is the ratio of buoyancy forces to viscous and thermal diffusion forces. Larger values of $Ra$ correspond to stronger buoyancy forces relative to viscous and thermal diffusion forces and drive more vigorous convection.For a given geometry and boundary conditions, there exists a critical value of $Ra$ below which no convection occurs, and the system undergoes a bifurcation from a stable conductive state to an unstable convective state.  


In this homework, we will solve the nonlinear system of equations for a range of $Ra$ values,  in a 2D square box with appropriate boundary conditions. But first we will do some preliminary work to simplify the problem and test our solvers. 

Do the following:
\begin{enumerate}
  \item (\textbf{10 pts}) \textbf{Rewrite the governing equations in stream-function/vorticity form}. In 2-D,  we can write the velocity field in terms of a scalar stream function $\psi$ 
  \begin{displaymath} 
    \mathbf{v} = \nabla \times (\psi \hat{\mathbf{k}}),
  \end{displaymath}
   where $\hat{\mathbf{k}}$ is the unit vector normal to the plane of flow and the incompressibility condition is automatically satisfied.  In addition, the vorticity is defined as 
   \begin{displaymath}
     \omega \hat{\mathbf{k}}= \nabla \times \mathbf{v}
   \end{displaymath}
    With these definitions, show that governing equations (Eqs. \ref{eq:temp}--\ref{eq:incomp}) can be rewritten as 
    \begin{align}
      \frac{\partial T}{\partial t} + \mathbf{v} \cdot \nabla T & = \frac{1}{Ra}\mathbf{\nabla}^2 T  \label{eq:temp2}\\
      -\nabla^2\omega & = \frac{\partial T}{\partial x} \label{eq:vort} \\
      -\nabla^2 \psi & = \omega \label{eq:stream}
    \end{align}
  \item (\textbf{10 pts}) \textbf{Explore fieldsplit preconditioners for solving the  biharmonic equation in PETSc and Firedrake}. I have provided two codes that solve just the biharmonic equation (Eqs. \eqref{eq:vort} and \eqref{eq:stream}) with a manufactured solution. The first code is written in C using the PETSc library (\texttt{biharm.c}) for finite differences on a \texttt{DMDA} and the second code is written in Python using the Firedrake library (\texttt{biharm.py}) for finite elements. Both codes have options for monolithic and  fieldsplit preconditioners to solve the coupled system of equations.  The Petsc C code includes three options files for different solver options. 
  \begin{enumerate}
    \item Compare convergence rates and run times for the three options files and the c code. You can use 
    \begin{verbatim}   
      ./biharm -options_file <file> -log_view | grep SNESSolve
    \end{verbatim}
    to run the code and get detailed timing information on the SNES solver (and other functions) 
    \item Repeat the exercise for the firedrake code using the same set of options.
    \item pick one solver set and compare the output log files for the two codes and determine where the principle differences in run time are coming from. 
  \end{enumerate}
  \item (\textbf{5 pts}) Write a new firedrake code based on \texttt{biharm.py} to solve the biharmonic equation (Eqs. \eqref{eq:temp2}--\eqref{eq:stream}) but with a RHS function $f(x) = \frac{\partial T}{\partial x}$ and \
  \begin{equation}
    T(x,y) = (1 - y) + A\cos(\pi x) \label{eq:Tinit}
  \end{equation}
  with $A=0.1$. (hint: if \texttt{T = Function(V)} is a firedrake \texttt{ufl} function, then \texttt{T.dx(0)} evaluates symbolically as $\frac{\partial T}{\partial x}$). Just provide a paraview figure showing the temperature field, with contours of the stream function and the velocity $v=\mathbb{\nabla}\times\psi\hat{\mathbf{k}}$ superposed.
  \item (\textbf{20 pts}) Write a firedrake program \texttt{convection.py} that solves the full system of equations as an implicit DAE problem using \texttt{firedrake-ts}.  You can start with your temperature driven biharmonic solver and the \texttt{firedrake-ts} code \texttt{heat.py} and combine them.  The specific problem is to solve Eqs. (\ref{eq:temp2}-\ref{eq:stream}) with Boundary conditions
  \begin{eqnarray*}
    T(x,0,t) = 1 & T(x,1,t) = 0\\
    \frac{\partial T}{\partial x}(0,y,t) = 0 & \frac{\partial T}{\partial x}(1,y,t)  = 0\\
    \omega = 0 & \psi = 0\quad\mathrm{on}\, \partial\Omega \\
  \end{eqnarray*}
  and initial condition Eq. \eqref{eq:Tinit} as a monolithic DAE $F(u,\dot(u),t)=0$.  There is no analytic solution to this problem, but you should calculate the Nusselt Number 
  \begin{displaymath}
    Nu = -\frac{\int_0^1 \partial_y T(x, y=1)}{\int_0^1 \partial_y T(x, y=0)}
  \end{displaymath}
  which is a measure of the convective heat flux to the conductive heat flux (for no-convection $Nu=1$)
  \begin{enumerate}
    \item Start with $Ra=10^2$ on a $64\times64$ Q1 mesh, use a bdf-2 integrator and a monolithic direct solver with mumps and solve to a maximum time $t=100000$.  Show that for this Ra, the steady state is non-convective.
    \item Repeat the exercise for $Ra=10^4,10^5,10^6$
    \item For $Ra=10^4$, plot $Nu$ vs mesh size for $N\times N$ meshes with $N=16,32,64,128$ and compare to the results of Blankenbach benchmark 1 which predicts $Nu=4.884$.  Discuss.
   \end{enumerate}
   \item \textbf{Extra credit themes and variations:} (not recommended, just for fun...better to work on your project)
   \begin{enumerate}
    \item Do full convergence study for $Ra=10^4,10^5,10^6$
    \item Rewrite the problem as a Petsc c code for finite difference on a DMDA using TS (with fieldsplit block precondtioners)
    \item Use Finite Difference in Time/Finite elements in space  with a 2nd order Crank-Nicholson discrization and write the problem as a SNES for each time step and use plain firedrake, but allow block pre-conditioners
   \end{enumerate}
  \end{enumerate}

  
  \end{document}